\section*{TD 1 --- Démarche et généralisation}
\addcontentsline{toc}{section}{TD 1 --- Démarche et généralisation}

\subsection*{Exercice 1 --- Vocabulaire \hfill \normalfont\footnotesize (5 min)}
Pour chaque situation, donner $X$ et $y$, puis le type de tâche.
\begin{enumerate}
  \item Prédire la consommation électrique d'un bâtiment à partir de l'heure,
        du jour et de la température.
  \item Regrouper des clients selon leurs achats, sans catégorie connue.
  \item Décider si un dossier de crédit est risqué, à partir de revenus et
        d'historiques de remboursement.
\end{enumerate}

\subsection*{Exercice 2 --- Risque empirique \hfill \normalfont\footnotesize (10 min)}
On dispose de quatre observations :

\begin{center}
\small
\begin{tabular}{@{}lrrrr@{}}
\toprule
$i$ & 1 & 2 & 3 & 4 \\
\midrule
$y_i$ & 100 & 150 & 80 & 200 \\
$\hat y_i$ & 90 & 120 & 100 & 180 \\
\bottomrule
\end{tabular}
\end{center}

\begin{enumerate}
  \item Calculer la perte quadratique moyenne $\hat R_4$ et la RMSE.
  \item Calculer la perte absolue moyenne.
  \item Une erreur de 30 pèse combien de fois plus qu'une erreur de 10 dans
        chaque perte ?
\end{enumerate}

\subsection*{Exercice 3 --- Diagnostiquer \hfill \normalfont\footnotesize (8 min)}
Trois modèles, trois bilans. Pour chacun : sous-apprentissage, bon compromis ou
surapprentissage ? Justifier en une phrase.

\begin{center}
\small
\begin{tabular}{@{}lccc@{}}
\toprule
& \textbf{Erreur entraînement} & \textbf{Erreur test} & \textbf{Diagnostic} \\
\midrule
A & 0,32 & 0,34 & \\
B & 0,02 & 0,41 & \\
C & 0,08 & 0,09 & \\
\bottomrule
\end{tabular}
\end{center}

\subsection*{Exercice 4 --- Choisir le découpage \hfill \normalfont\footnotesize (8 min)}
Quel découpage (aléatoire, temporel, groupé, stratifié) imposez-vous ?
\begin{enumerate}
  \item Prédire les ventes du mois suivant à partir de 3 ans d'historique.
  \item Détecter une maladie rare : 4 \% de cas positifs.
  \item Classer des élèves dont certains ont plusieurs relevés annuels.
\end{enumerate}

\subsection*{Exercice 5 --- Fuite ou pas ? \hfill \normalfont\footnotesize (8 min)}
Un groupe prédit la réussite d'étudiants en octobre. Son notebook contient :
\begin{itemize}
  \item « moyenne du premier semestre » ;
  \item « nombre d'absences » ;
  \item « a été convoqué en tutorat en juin » ;
  \item une standardisation calculée sur l'ensemble des données avant le split.
\end{itemize}
Identifier les deux problèmes et proposer une correction pour chacun.


\subsection*{Exercice 6 --- Exactitude sur huit dossiers \hfill \normalfont\footnotesize (6 min)}
Un modèle classe 8 dossiers : 6 sont correctement prédits.
\begin{enumerate}
  \item Calculer l'exactitude.
  \item Ce chiffre suffit-il si 7 dossiers sur 8 sont négatifs ? Justifier.
  \item Quelle information manque pour conclure ?
\end{enumerate}

\subsection*{Exercice 7 --- Choisir la perte \hfill \normalfont\footnotesize (6 min)}
On prédit le \textbf{nombre d'incidents} mensuels d'un serveur (comptage).
\begin{enumerate}
  \item Pourquoi la perte quadratique est-elle peu adaptée ?
  \item Quelle perte privilégier pour un comptage ?
  \item Que se passe-t-il si l'on prédit une valeur négative ?
\end{enumerate}

\subsection*{Exercice 8 --- Le score trop beau \hfill \normalfont\footnotesize (5 min)}
Un groupe annonce 100 \% d'exactitude sur un problème où les données contiennent
du bruit. Rédiger les trois vérifications à demander avant de le croire.

\subsection*{Exercice 9 --- 60/20/20 ou validation croisée ? \hfill \normalfont\footnotesize (6 min)}
Donner un avantage et un inconvénient de chaque approche, puis indiquer laquelle
choisir avec 120 observations.

\newpage
